% Part: lambda-calculus
% Chapter: lambda-definability
% Section: introduction

\documentclass[../../../include/open-logic-section]{subfiles}

\begin{document}

\olfileid{lam}{ldf}{int}
\olsection{引言}

乍看之下，λ-演算仅仅是一个极其抽象的表达式演算，其中的表达式表示函数以及函数对其他函数的应用。λ-演算的语法丝毫未暗示这些是作用于特定种类对象的函数；特别是，其语法完全没有提及自然数。它的基本操作——应用与λ抽象——是适用于任意函数的操作，而不仅限于自然数上的函数。

然而，只要稍加巧思，便有可能在λ-演算中定义算术函数，即自然数上的函数。为此，对于每个自然数~$n \in \Nat$，我们定义一个特殊的λ-项~$\num n$，即~$n$ 的\emph{丘奇数}。（丘奇数得名于 邱奇。）

\begin{defn}
  若 $n \in \Nat$，则对应的\emph{丘奇数} $\num{n}$ 表示~$n$：
  \[
    \num{n} \ident \lambd[fx][f^n(x)]
  \]
  这里，$f^n(x)$ 表示将 $f$ 对~$x$ 应用 $n$ 次所得的结果。
  例如，$\num{0}$ 是 $\lambd[fx][x]$，而 $\num{3}$ 是
  $\lambd[fx][f(f(f\,x))]$。
\end{defn}
  
丘奇数 $\num n$ 被编码为一个λ-项，该项表示一个接受两个参数 $f$ 和 $x$ 并返回 $f^n(x)$ 的函数。丘奇数显然均处于范式。

当然，只有能在λ-演算中用自然数的这种表示进行计算，它才是有用的。在λ-演算中用丘奇数进行计算，是指将一个λ-项~$F$ 应用于这样一个丘奇数，并将组合项~$F\, \num n$ 归约为范式。如果该组合项总能归约为范式，且该范式总是一个丘奇数~$\num m$，我们就可以认为该计算的输出为数~$m$。于是，我们可以将~$F$ 视作定义了一个函数 $f\colon \Nat \to \Nat$，即满足 $f(n) = m$ 当且仅当 $F\, \num n \red \num m$ 的函数。根据 邱奇--罗瑟 性质，范式若存在则唯一。因此，如果 $F\, \num n \red \num m$，$F \, \num n$ 便不可能归约到其他任何范式，尤其不可能归约到另一个丘奇数。

反之，给定一个函数 $f\colon \Nat \to \Nat$，我们可以问是否存在一个λ-项 $F$ 以这种方式定义~$f$。若存在这样的项，我们便称 $F$ \emph{λ-定义}了~$f$，并称 $f$ 是λ-可定义的。我们可将此概念推广至多元函数与部分函数。

\begin{defn}
设 $f\colon \Nat^k \to \Nat$。我们称一个λ-项~$F$ \emph{λ-定义}了 $f$，如果对所有 
$n_0$, \dots,~$n_{k-1}$，
\[
F \, \num{n_0} \, \num{n_1} \dots \num{n_{k-1}} \red \num{f(n_0, n_1, \dots,
  n_{k-1})}
\]
若 $f(n_0, \dots, n_{k-1})$ 有定义，否则 $F \, \num{n_0} \,
\num{n_1} \dots \num{n_{k-1}}$ 无范式。
\end{defn}

常数函数是很简单的例子。项 $C_k \ident
\lambd[x][\num{k}]$ λ-定义了函数 $c_k\colon \Nat \to
\Nat$，使得 $c(n) = k$。因为对任意~$n$，有 $C_k \, \num n \ident
(\lambd[x][\num{k}])\num n \redone \num{k}$。恒等函数由 $\lambd[x][x]$ λ-定义。当然，更复杂的函数更难定义，通常需要极大的巧思。因此，每个可计算函数都是λ-可定义的，这或许令人惊讶。反之亦然：如果一个函数是λ-可定义的，则它是可计算的。

\end{document}